\section{Conclusion and Discussion}
\label{sec:conclusion}
In this paper, we introduced an algorithm for training gender-neutral word embedding. Our method is general and can be applied in any language as long as a list of gender definitional words is provided as seed words (e.g., gender pronouns). Future directions include extending the proposed approach to model other properties of words such as sentiment and generalizing our analysis beyond binary gender. 

\section*{Acknowledgement}
This work was supported by National Science Foundation Grant IIS-1760523. We would like to thank James Zou, Adam Kalai, Tolga Bolukbasi, and Venkatesh Saligrama for helpful discussion, and anonymous reviewers for their feedback.

%The analysis can be easily extended to languages without grammatical gender such as Chinese, English and Japanese. 
%The method can be generalized to accommodate other demographic biases or semantic features of word embeddings.
%Any binary stereotypes or features, e.g. ``old or young'' and ``racial discriminatory or not'', can be handled with a small set of feature-specific words and a predefined feature direction from the word embedding. 
%Also GN-GloVe can be generalized to other domains other than gender.
% For example, it can be used to learn sentimental embeddings for the sentiment analysis tasks. 
%In the future, we consider extending our approach to model other properties of words, such as sentiment.  
%other domains, such as encoding sentiment. For example, current word embeddings represent ``good'' and ``bad'' using similar vectors as they often occur under similar context. 
%With our method, the sentiment meaning of such words will be far away from each other in  specifically defined dimensions, which can be helpful to the sentiment analysis tasks. 
%Besides the binary gender cases, we can generalize our model to multi-class cases by using $k$ dimensions for the protected feature.
%, where each dimension is a binary value for that specific label.